\begin{thebibliography}{9}

\bibitem{wastlund}
J.~W\"astlund, \emph{Summing inverse squares by euclidean geometry},
Chalmers University of Technology, Gothenburg, 8 December 2010.
Unpublished note.
\url{https://www.math.chalmers.se/~wastlund/Cosmic.pdf}

\bibitem{threeblueonebrown}
3Blue1Brown, \emph{Why is $\pi$ here? And why is it squared? A geometric answer to the
Basel problem}, YouTube, 3 March 2018.
\url{https://www.youtube.com/watch?v=d-o3eB9sfls}

\bibitem{kalman}
D.~Kalman, \emph{Six ways to sum a series}, College Mathematics Journal 24 (5), 402--421, 1993.

\bibitem{matsuoka}
Y.~Matsuoka, \emph{An elementary proof of the formula $\sum 1/k^{2}=\pi^{2}/6$},
American Mathematical Monthly 68 (5), 1961.

\bibitem{komornik}
V.~Komornik, \emph{A ``luminous'' solution to the Basel problem},
Mathematics Magazine 95 (4), 333--337, 2022.

\bibitem{wikipedia}
{Wikipedia contributors}, \emph{Basel problem}, section ``Geometric proof'', which reproduces
the construction of \cite{wastlund}.
\url{https://en.wikipedia.org/wiki/Basel_problem}, accessed 20 November 2026.

\end{thebibliography}
